\cgeentry{cge_elasticities}{Published elasticities, with their sources}

\cgesec{Description}

\cmd{cge_elasticities} lists, loads, or turns into values for the options of
\manref{cge_calibrate}, the entries of the package's elasticity library. Each
entry is a published estimate, with its source, the table and page where it is
printed, its definition and the check made when it was read from the source.
Elasticities have no defaults in \manref{cge_calibrate}; the library documents
where values come from, and nothing is applied without a command. Entries of
the user's own, from a file in the library's layout, may be added.

\cgesec{Quick start}

List the Armington elasticities of Togo, Sierra Leone and Botswana
\begin{stcode}
cge_elasticities sigmaq, country(TGO SLE BWA)
\end{stcode}
Load the income elasticities of Togo, with their definitions and sources, into
memory
\begin{stcode}
cge_elasticities incelas, country(TGO) clear
\end{stcode}
After \manref{cge_map}, one Armington elasticity per commodity, for
\cmd{sigmaq()}
\begin{stcode}
cge_elasticities sigmaq, country(TGO) values
local sq `r(values)'
\end{stcode}
Income elasticities by commodity, with the consumption category of each
commodity
\begin{stcode}
cge_elasticities incelas, country(TGO) values map(com_agr:food com_hlt:health) others(other)
\end{stcode}
Add the entries of a file to the library and list those of one source,
\cmd{MY2026}
\begin{stcode}
cge_elasticities using myvalues.csv, source(MY2026)
\end{stcode}

\cgesec{Syntax}

\begin{stsyntax}
\cmdb{cge_elasticities} [\argu{parameter}] [\texttt{using} \argu{filename}]
[\texttt{,} \argu{options}]
\end{stsyntax}

\argu{parameter} is the \manref{cge_calibrate} option that a value is for:
\cmd{sigmaq}, \cmd{sigmat}, \cmd{sigmava}, \cmd{incelas}, \cmd{frisch},
\cmd{wagecurve} or \cmd{mortelas} in the package's library.

\begin{opttable}
\optgroup{Selection}
\optarg{c}{ountry}{codes} & entries for these countries or groups\\
\optarg{cat}{egory}{names} & entries in these categories\\
\optarg{sou}{rce}{ids} & entries from these sources\\
\optarg{}{id}{ids} & these entries\\
\optgroup{Output}
\opt{}{clear} & load the selected entries into memory\\
\opt{val}{ues} & return values for \manref{cge_calibrate} in \res{r(values)}\\
\optarg{}{map}{spec} & with \cmd{values}, categories of commodities or activities\\
\optarg{oth}{ers}{category} & with \cmd{values}, the category of those not in \cmd{map()}\\
\end{opttable}

\cgesec{Options}

\begin{optdesc}
\item[\optarg{c}{ountry}{codes}, \optarg{cat}{egory}{names},
\optarg{sou}{rce}{ids} and \optarg{}{id}{ids}] select entries. Each takes one
or more values, and an entry is selected when it matches every option given,
and \argu{parameter} if it is given. Countries are ISO 3166 alpha-3 codes, for
example \cmd{TGO}, or the groups \cmd{AFR} (Africa), \cmd{WLD} (international)
and \cmd{DEV} (developing countries). The case of the values does not matter.
If no entry is selected, the command stops with error 111.

\item[\cmd{using} \argu{filename}] adds the entries of a CSV file in the
library's layout, with at least the columns \cmd{id}, \cmd{parameter},
\cmd{country}, \cmd{category}, \cmd{value} and \cmd{source}, for example values
that the user is licensed to use or has estimated. Other columns of the layout
are empty when absent. Entry identifiers must be unique across the library and
the file, and every entry needs a numeric value.

\item[\opt{}{clear}] loads the selected entries into memory, replacing the
data, with every column of the library and the authors, year, title,
publication and address of each source. \cmd{value}, \cmd{se}, \cmd{lower} and
\cmd{upper} are numeric; the other variables are strings. \cmd{clear} may not
be combined with \cmd{values}.

\item[\opt{val}{ues}] returns values for \manref{cge_calibrate} in
\res{r(values)} and displays the entries used, with their definitions, sources
and notes. For \cmd{sigmaq}, \cmd{sigmat} and \cmd{incelas} it returns one value
per commodity, and for \cmd{sigmava} one per activity, in the order of
\manref{cge_map}, which must have been run. Each commodity or activity takes
the selected entry of its category, given by \cmd{map()} or \cmd{others()}, or
the single selected entry if it has no category; if no entry or more than one
matches, the command stops and names the entries. For \cmd{frisch},
\cmd{wagecurve}, \cmd{mortelas} and \cmd{sickelas} it returns one number, from
the single selected entry; \cmd{wagecurve()} of \manref{cge_calibrate} takes it
with the factor, as in
\texttt{wagecurve(lab:\textasciigrave r(values)\textquotesingle)}. \argu{parameter} is
required.

\item[\optarg{}{map}{spec}] gives categories to commodities or activities, as
\argu{element}\cmd{:}\argu{category} pairs, for example
\cmd{map(com_agr:food com_hlt:health)}. It needs \cmd{values}.

\item[\optarg{oth}{ers}{category}] gives a category to the commodities or
activities not listed in \cmd{map()}. It needs \cmd{values}.
\end{optdesc}

\cgesec{Remarks and examples}

The library holds 531 entries (\chapref{chap:elasticities}):

\begin{cgetable}{>{\raggedright\arraybackslash}p{0.13\linewidth}>{\raggedright\arraybackslash}p{0.08\linewidth}>{\raggedright\arraybackslash}p{0.41\linewidth}>{\raggedright\arraybackslash}p{0.30\linewidth}}
\toprule
Parameter & Entries & Coverage & Source\\
\midrule
\cmd{sigmaq} & 51 & aggregate Armington elasticities, African countries & Devarajan, Go and Robinson (2023)\\
\cmd{sigmat} & 51 & aggregate CET elasticities, the same countries & Devarajan, Go and Robinson (2023)\\
\cmd{incelas} & 423 & income elasticities for nine consumption categories, 47 African countries & Muhammad et al.\ (2011)\\
\cmd{frisch} & 1 & the Frisch parameter, 144 countries pooled & Muhammad et al.\ (2011)\\
\cmd{sigmava} & 1 & substitution between capital and labour, 30 African countries pooled & Hosni and Ouakil (2024)\\
\cmd{wagecurve} & 2 & international, and South Africa & Blanchflower and Oswald (1995); Kingdon and Knight (2006)\\
\cmd{mortelas} & 2 & under-five and maternal mortality with respect to government health expenditure & Bokhari, Gai and Gottret (2007)\\
\bottomrule
\end{cgetable}

The categories of the income elasticities are \cmd{food} (with beverages and
tobacco), \cmd{clothing}, \cmd{housing}, \cmd{furnishings}, \cmd{health},
\cmd{transport} (with communication), \cmd{recreation}, \cmd{education} and
\cmd{other}; the trade elasticities have the category \cmd{aggregate}. The
mortality elasticities are for reference: their outcome and driver differ from
those of \cmd{mortelas()} of \manref{cge_calibrate}.

Read the definitions before using a value. Aggregate trade elasticities are
lower than sector ones, so they set a level rather than every sector's value
(Devarajan, Go and Robinson 2023). The income elasticities come from a
cross-country demand system for 2005, not from a household survey of the
country. Two values follow the package's conventions rather than the sources:
the wage-curve elasticity is the absolute value of the published negative
elasticity, and the Frisch parameter is the inverse of the published income
flexibility.

\cgesubsec{Example 1: listing and loading entries}

Without \cmd{clear} or \cmd{values}, the selected entries are listed, up to 60,
with the location of each value in its source:
\stlog[\scriptsize]{elas_list}
With \cmd{clear}, the entries are loaded with all their details:
\stlog{elas_load}
The library has no standard errors for these income elasticities.

\cgesubsec{Example 2: values for calibration}

The values are given for the commodities and activities of a mapped model,
here the synthetic SAM:
\stlog{elas_setup}
\stlog{elas_values}
The one Armington elasticity of Togo is given to every commodity. The income
elasticities need a category for each commodity: agriculture takes the
elasticity of food, and the other commodities that of other goods and services.
The values are then options of \manref{cge_calibrate}:
\stlog{elas_calibrate}
With two countries selected, each commodity matches two entries, and the
command stops:
\stlog{elas_ambiguous}
The results name the entry used for each element:
\stlog{elas_return}

\cgesubsec{Example 3: a file of the user's own entries}

A file of additional entries has one row per entry and at least the columns
\cmd{id}, \cmd{parameter}, \cmd{country}, \cmd{category}, \cmd{value} and
\cmd{source}. The library's other columns, \cmd{country_name}, \cmd{se},
\cmd{lower}, \cmd{upper}, \cmd{published}, \cmd{definition}, \cmd{method},
\cmd{observations}, \cmd{period}, \cmd{location}, \cmd{check} and \cmd{notes},
may be added. Here a file gives two sector Armington elasticities, with their
own categories:
\stlog{elas_user}
The file has no \cmd{location} column, so none is displayed after the source.

\cgesec{Stored results}

\cmd{cge_elasticities} stores the following in \cmd{r()}:
\begin{storedres}
\resgroup{Scalars}
\res{r(N)} & number of entries selected, or of values returned\\
\resgroup{Macros}
\res{r(ids)} & entries selected, or the entry used for each value\\
\res{r(values)} & values for \manref{cge_calibrate} (with \cmd{values})\\
\res{r(parameter)} & parameter (with \cmd{values})\\
\end{storedres}

\cgesec{References}

\begin{list}{}{\setlength{\leftmargin}{2em}\setlength{\itemindent}{-2em}%
  \setlength{\itemsep}{2pt}\setlength{\parsep}{0pt}}
\item Blanchflower, D. G., and A. J. Oswald. 1995. An introduction to the wage
curve. \textit{Journal of Economic Perspectives} 9(3): 153--167.
\item Bokhari, F. A. S., Y. Gai, and P. Gottret. 2007. Government health
expenditures and health outcomes. \textit{Health Economics} 16(3): 257--273.
\item Devarajan, S., D. S. Go, and S. Robinson. 2023. Trade elasticities in
aggregate models: Estimates for 191 countries. World Bank Policy Research
Working Paper 10490.
\item Hosni, M., and H. Ouakil. 2024. Estimation de l'\'elasticit\'e de
substitution entre le capital et le travail de l'\'economie africaine~: une
analyse en donn\'ees de panel. \textit{Revue AME} 6(3): 247--269.
\item Kingdon, G. G., and J. Knight. 2006. How flexible are wages in response
to local unemployment in South Africa? \textit{ILR Review} 59(3): 471--495.
\item Muhammad, A., J. L. Seale Jr., B. Meade, and A. Regmi. 2011.
\textit{International Evidence on Food Consumption Patterns: An Update Using
2005 International Comparison Program Data}. Technical Bulletin TB-1929,
Economic Research Service, U.S. Department of Agriculture.
\end{list}

\cgesec{Also see}

\begin{alsosee}
\seeentry{cge_calibrate}{Calibrate the model to the SAM}
\seeentry{cge_sensitivity}{Sensitivity of a scenario's results to the elasticities}
\seechapter{chap:elasticities}{The elasticity library}
\end{alsosee}
